\ifdefined\gkver\else\def\gkver{student}\fi
\documentclass[\gkver]{gaokaozhenti}
\newcommand{\ccnum}[1]{{\fontspec{Noto Serif CJK SC}#1}}
\begin{document}

\chapter{2014年全国课标Ⅰ}
%% paperType: 理综物理部分

\begin{enumerate}
\item
%% number: 14
%% paperName: 2014年全国课标Ⅰ第 14 题 6 分
%% typeId: 1
%% type: 单选题
%% chapter: 电磁感应
%% point: 电磁感应现象
%% method: 无
%% score: 6
%% degree: 650
%% duplicateId: 0
%% body:
在法拉第时代，下列验证“由磁产生电”设想的实验中，能观察到感应电流的是\xzanswer{D}
\fourchoices[answer=D]
{将绕在磁铁上的线圈与电流表组成一闭合回路，然后观察电流表的变化}
{在一通电线圈旁放置一连有电流表的线圈，然后观察电流表的变化}
{将一房间内的线圈两端与相邻房间的电流表相连。往线圈中插入条形磁铁后，再到相邻房间去观察电流表的变化}
{绕在同一铁环上的两个线圈，分别接电源和电流表，在给线圈通电或断电的瞬间，观察电流表的变化}
%% answer: D
%% memo:
\memoanswer{
法拉第发现的电磁感应定律并总结出五种情况下会产生感应电流，其核心就是穿过闭合线圈的磁通量发生变化，选项 A、B 中，绕在磁铁上面的线圈和通电线圈，线圈面积都没有发生变化，前者磁场强弱没有变化，后者通电线圈中若为恒定电流则产生恒定的磁场，也是磁场强弱不变，都会导致磁通量不变化，不会产生感应电流，A、B 错误。选项 C 中往线圈中插入条形磁铁导致磁通量发生变化，在这一瞬间会产生感应电流，但是过程短暂，等到插入后再到相邻房间去，过程已经结束，观察不到电流表的变化。C 错误。选项 D 中，线圈通电或断电瞬间，导致线圈产生的磁场变化，从而引起另一个线圈的磁通量变化产生感应电流，可以观察到电流表指针偏转，D 正确。
}

\item
%% number: 15
%% paperName: 2014年全国课标Ⅰ第 15 题 6 分
%% typeId: 1
%% type: 单选题
%% chapter: 磁场
%% point: 安培力
%% method: 无
%% score: 6
%% degree: 650
%% duplicateId: 0
%% body:
关于通电直导线在匀强磁场中所受的安培力，下列说法正确的是\xzanswer{B}
\fourchoices[answer=B]
{安培力的方向可以不垂直于直导线}
{安培力的方向总是垂直于磁场的方向}
{安培力的大小与通电直导线和磁场方向的夹角无关}
{将直导线从中点折成直角，安培力的大小一定变为原来的一半}
%% answer: B
%% memo:
\memoanswer{
根据匀强磁场对直线电流的作用力可知：安培力的方向与导线和磁场方向相互垂直，故 A 错误，B 正确；由安培力的大小的计算公式 $F=BIL\sin\alpha$，$\alpha$ 是磁场方向和电流方向的夹角，所以安培力大小与通电导线和磁场方向的夹角有关，C 错误；将直导线从中点折成直角，由于其余磁场方向的夹角未知，其合力不一定等于原来的一半，所以 D 错误。
}

\item
%% number: 16
%% paperName: 2014年全国课标Ⅰ第 16 题 6 分
%% typeId: 1
%% type: 单选题
%% chapter: 磁场
%% point: 带电粒子在磁场中的运动
%% method: 无
%% score: 6
%% degree: 650
%% duplicateId: 0
%% body:
如图，MN 为铝质薄平板，铝板上方和下方分别有垂直于图平面的匀强磁场（未画出）。一带电粒子从紧贴铝板上表面的 P 点垂直于铝板向上射出，从 Q 点穿越铝板后到达 PQ 的中点 O。已知粒子穿越铝板时，其动能损失一半，速度方向和电荷量不变。不计重力。铝板上方和下方的磁感应强度大小之比为\xzanswer{D}
\onepicture[align=right, w=5.0cm]{figs/16.png}
\fourchoices[answer=D]
{2}
{$\sqrt 2 $}
{1}
{$\frac{{\sqrt 2 }}{2}$}
%% answer: D
%% memo:
\memoanswer{
粒子垂直于磁场方向进入磁场，在洛伦兹力作用下做匀速圆周运动 $qvB=m\frac{v^{2}}{R}$ 得：$B=\frac{mv}{qR}$；又粒子的动能：$E_{k}=\frac{1}{2}mv^{2}$，故 $B=\frac{\sqrt{2mE_{k}}}{qR}$。由题知：$R_{1}=2R_{2}$，$E_{k1}=2E_{k2}$，代入得：$\frac{B_{1}}{B_{2}}=\frac{\sqrt{2}}{2}$，D 正确。
}

\item
%% number: 17
%% paperName: 2014年全国课标Ⅰ第 17 题 6 分
%% typeId: 1
%% type: 单选题
%% chapter: 运动和力
%% point: 牛顿第二定律
%% method: 无
%% score: 6
%% degree: 450
%% duplicateId: 0
%% body:
如图，用橡皮筋将一小球悬挂在小车的架子上，系统处于平衡状态。现使小车从静止开始向左加速，加速度从零开始逐渐增大到某一值，然后保持此值，小球稳定地偏离竖直方向某一角度（橡皮筋在弹性限度内）。与稳定在竖直时位置相比，小球的高度\xzanswer{A}
\onepicture[align=right, w=4.2cm]{figs/17.png}
\fourchoices[answer=A]
{一定升高}
{一定降低}
{保持不变}
{升高或降低由橡皮筋的劲度系数决定}
%% answer: A
%% memo:
\memoanswer{
小车静止时，橡皮筋弹力等于小球重力 $F_{1}=mg$，小球在悬点下竖直距离：$L_{1}=L_{0}+mg/k$；小车以一定加速度运动时，小球稳定地偏离竖直方向某一角度 $\theta$，此时橡皮筋弹力为 $F_{2}$，对于小球在竖直方向有：$F_{2}\cos\theta=mg$，则小球在悬点下竖直距离：$L_{2}=(L_{0}+mg/k\cos\theta)\cos\theta=L_{0}\cos\theta+mg/k$，故 $L_{1}>L_{2}$，即小球一定升高。A 正确。
}

\item
%% number: 18
%% paperName: 2014年全国课标Ⅰ第 18 题 6 分
%% typeId: 1
%% type: 单选题
%% chapter: 电磁感应
%% point: 电磁感应中图象
%% method: 无
%% score: 6
%% degree: 450
%% duplicateId: 0
%% body:
如图（a），线圈 ab、cd 绕在同一软铁芯上，在 ab 线圈中通以变化的电流，用示波器测得线圈 cd 间的电压如图（b）所示。已知线圈内部的磁场与流经线圈的电流成正比，则下列描述线圈 ab 中电流随时间变化关系的图中，可能正确的是\xzanswer{C}
\twopicture[labelA=2014全国课标Ⅰ18a, labelB=2014全国课标Ⅰ18b, w=5.0cm]
{figs/18a.png}
{figs/18b.png}
\fourchoices[answer=C, ispicture=true, h=2.7cm]
{figs/18c.png}
{figs/18d.png}
{figs/18e.png}
{figs/18f.png}
%% answer: C
%% memo:
\memoanswer{
根据法拉第电磁感应定律，感应电动势即 cd 线圈中的电压 $U_{cd}=\frac{\Delta\Phi}{\Delta t}=\frac{\Delta B}{\Delta t}S$，由于磁场是线圈 ab 中的电流产生的，且线圈内部的磁场与流经线圈的电流成正比 $\Delta B\propto\Delta I$，故 $U_{cd}\propto\frac{\Delta I}{\Delta t}$。在 $i-t$ 图像中图线的斜率即为 $U_{cd}$，对照图（b）和选项图得 C 正确。
}

\item
%% number: 19
%% paperName: 2014年全国课标Ⅰ第 19 题 6 分
%% typeId: 2
%% type: 多选题
%% chapter: 曲线运动
%% point: 天体运动
%% method: 无
%% score: 6
%% degree: 450
%% duplicateId: 0
%% body:
太阳系各行星几乎在同一平面内沿同一方向绕太阳做圆周运动。当地球恰好运行到某地外行星和太阳之间，且三者几乎排成一条直线的现象，天文学称为“行星冲日”。据报道，2014年各行星冲日时间分别是：1月6日木星冲日；4月9日火星冲日；5月11日土星冲日；8月29日海王星冲日；10月8日天王星冲日。已知地球及各地外行星绕太阳运动的轨道半径如下表所示，则下列判断正确的是\xzanswer{BD}

\begin{center}
\begin{tabular}{|c|c|c|c|c|c|c|}
\hline
 & 地球 & 火星 & 木星 & 土星 & 天王星 & 海王星 \\ \hline
轨道半径（$\UAU$） & 1.0 & 1.5 & 5.2 & 9.5 & 19 & 30 \\ \hline
\end{tabular}
\end{center}

\fourchoices[answer=BD]
{各地外行星每年都会出现冲日现象}
{在2015年内一定会出现木星冲日}
{天王星相邻两次冲日的时间间隔为土星的一半}
{地外行星中，海王星相邻两次冲日的时间间隔最短}
%% answer: BD
%% memo:
\memoanswer{
相邻两次行星冲日的时间间隔就是地球比该行星多运动一周的时间，根据开普勒第三定律 $\frac{R^{3}}{T^{2}}=k$ 得：$T=\sqrt{R^{3}/k}$，相邻两次行星冲日的时间间隔 $t=\frac{2\pi}{\omega_{\text{地}}-\omega_{\text{行}}}=\frac{2\pi}{\frac{2\pi}{T_{\text{地}}}-\frac{2\pi}{T_{\text{行}}}}=\frac{T_{\text{地}}}{1-\frac{T_{\text{地}}}{T_{\text{行}}}}>T_{\text{地}}$，即相邻两次行星冲日的时间间隔大于 1 年，A 错误；根据木星轨道半径是地球的 5.2 倍，木星周期大于 11 年，小于 12 年，所以木星冲日的时间间隔大于 $\frac{12}{11}$ 年，小于 1.1 年。由于今年的冲日时间是 1 月 6 日，所以下次木星冲日在 2015 年，B 正确；由于木星相邻两次冲日的时间间隔的一半还不到一年，而所有行星冲日的相邻间隔都超过 1 年，所以天王星相邻两次的冲日的时间不可能是土星的一半，C 错误。根据海王星的轨道半径最大，周期最大，可判断海王星的时间间隔最短，D 正确。
}

\item
%% number: 20
%% paperName: 2014年全国课标Ⅰ第 20 题 6 分
%% typeId: 2
%% type: 多选题
%% chapter: 曲线运动
%% point: 圆周运动的应用
%% method: 临界
%% score: 6
%% degree: 450
%% duplicateId: 0
%% body:
如图，两个质量均为 $m$ 的小木块 a 和 b（可视为质点）放在水平圆盘上，a 与转轴 OO$'$ 的距离为 $l$，b 与转轴的距离为 $2l$。木块与圆盘的最大静摩擦力为木块所受重力的 $k$ 倍，重力加速度大小为 $g$。若圆盘从静止开始绕转轴缓慢地加速转动，用 $\omega$ 表示圆盘转动的角速度，下列说法正确的是\xzanswer{AC}
\onepicture[align=right, w=4.2cm]{figs/20.png}
\fourchoices[answer=AC]
{b 一定比 a 先开始滑动}
{a、b 所受的摩擦力始终相等}
{$\omega=\sqrt{\frac{kg}{2l}}$ 是 b 开始滑动的临界角速度}
{当 $\omega=\sqrt{\frac{2kg}{3l}}$ 时，a 所受摩擦力的大小为 $kmg$}
%% answer: AC
%% memo:
\memoanswer{
小木块都随圆盘做匀速圆周运动时，在发生相对滑动之前，角速度相等，静摩擦力提供向心力，即 $F_{\text{静}}=m\omega^{2}r$，由于 $r_{a}=l$，$r_{b}=2l$，所以发生相对滑动前木块 b 的静摩擦力大，B 错误；随着角速度的增大，当静摩擦力等于最大静摩擦力，木块将开始滑动，即 $kmg=m\omega^{2}r$，$\omega=\sqrt{\frac{kg}{r}}$，带入两个木块的半径，小木块 a 开始滑动时的角速度大于木块 b 开始滑动时的角速度。A、C 正确；又 $\omega=\sqrt{\frac{2kg}{3l}}<\sqrt{\frac{kg}{l}}$，所以木块 a 未达到临界状态，摩擦力还没有达到最大静摩擦力，所以选项 D 错误。
}

\item
%% number: 21
%% paperName: 2014年全国课标Ⅰ第 21 题 6 分
%% typeId: 2
%% type: 多选题
%% chapter: 电场
%% point: 等势面
%% method: 无
%% score: 6
%% degree: 450
%% duplicateId: 0
%% body:
如图，在正点电荷 Q 的电场中有 M、N、P、F 四点，M、N、P 为直角三角形的三个顶点，F 为 MN 的中点，$\angle M=30^{\circ}$。M、N、P、F 四点处的电势分别用 $\varphi_{M}$、$\varphi_{N}$、$\varphi_{P}$、$\varphi_{F}$ 表示，已知 $\varphi_{M}=\varphi_{N}$，$\varphi_{P}=\varphi_{F}$，点电荷 Q 在 M、N、P 三点所在平面内，则\xzanswer{AD}
\onepicture[align=right, w=4.5cm]{figs/21.png}
\fourchoices[answer=AD]
{点电荷 Q 一定在 MP 的连线上}
{连接 PF 的线段一定在同一等势面上}
{将正试探电荷从 P 点搬运到 N 点，电场力做负功}
{$\varphi_{P}$ 大于 $\varphi_{M}$}
%% answer: AD
%% memo:
\memoanswer{
点电荷电场的等势面是以点电荷为圆心的一族同心圆，所以任意两个等势点连线的垂直平分线都指向点电荷，如图所示，MN 的垂直平分线和 PE 的垂直平分线相交于一点 E，E 点即点电荷所在的位置，根据几何关系可得 E 点在 MP 边。即点电荷一定在 MP 连线上，选项 A 正确。点电荷的等势面是以点电荷为圆心的同心圆，所以等势面不是直线而是球面，B 错误；正试探电荷从 P 点移动到 N 点，远离场源正电荷，电场力做正功，C 错误。根据几何关系可得 EF < EM，距离场源正电荷越远电势越低，所以 P 点电势大于 M 点电势，D 正确。

\onepicture[align=right, w=5.0cm]{figs/21_an.png}
}

\item
%% number: 22
%% paperName: 2014年全国课标Ⅰ第 22 题 6 分
%% typeId: 4
%% type: 实验题
%% chapter: 运动和力
%% point: 验证牛顿第二定律
%% method: 无
%% score: 6
%% degree: 450
%% duplicateId: 0
%% body:
某同学利用图（a）所示实验装置及数字化信息系统获得了小车加速度 $a$ 与钩码的质量 $m$ 的对应关系图，如图（b）所示。实验中小车（含发射器）的质量为 $200\Ug$，实验时选择了不可伸长的轻质细绳和轻定滑轮，小车的加速度由位移传感器及与之相连的计算机得到。回答下列问题：
\begin{enumerate}
	\item
	根据该同学的结果，小车的加速度与钩码的质量成\_\_\_\_\_\_\_（填“线性”或“非线性”）关系。
	\item
	由图（b）可知，$a-m$ 图线不经过原点，可能的原因是\_\_\_\_\_\_\_\_\_\_\_\_\_\_。
	\item
	若利用本实验装置来验证“在小车质量不变的情况下，小车的加速度与作用力成正比”的结论，并直接以钩码所受重力 $mg$ 作为小车受到的合外力，则实验中应采取的改进措施是\_\_\_\_\_\_\_\_\_\_\_\_\_\_\_，钩码的质量应满足的条件是\_\_\_\_\_\_\_\_\_\_\_\_\_\_\_\_\_\_\_。
\end{enumerate}
\twopicture[labelA=2014全国课标Ⅰ22a, labelB=2014全国课标Ⅰ22b, hA=2.2cm, hB=4.5cm, align=right]
{figs/22a.png}
{figs/22b.png}
%% answer:
\jdanswer{
\begin{enumerate}
	\item
	非线性

	\item
	存在摩擦力

	\item
	调节轨道的倾斜度以平衡摩擦力，远小于小车的质量

\end{enumerate}
}
%% memo:
\memoanswer{
（1）根据该同学描绘的加速度和钩码质量的图像是一条曲线而不是直线，所以是非线性关系。

（2）根据图（b）可知当钩码有一定质量，即细线有一定拉力时，小车加速度仍等于 0，说明小车合力等于 0，所以可能除拉力外小车还受到摩擦力作用。

（3）实验改进部分由两个要求，第一个就是图像不过原点，需要平衡摩擦力，所以调整轨道倾斜度，第二个就是图像是曲线，因为小车的合力即细线拉力并不等于钩码的重力，而加速度 $a=\frac{mg}{m_{\text{车}}+m}$，所以要近似为直线，就要求钩码的质量远小于小车的质量，$k=\frac{g}{m_{\text{车}}+m}$ 近似为常量。
}

\item
%% number: 23
%% paperName: 2014年全国课标Ⅰ第 23 题 9 分
%% typeId: 4
%% type: 实验题
%% chapter: 电路
%% point: 测电动势与内阻
%% method: 无
%% score: 9
%% degree: 450
%% duplicateId: 0
%% body:
利用如图（a）所示电路，可以测量电源的电动势和内阻，所用的实验器材有：

待测电源，电阻箱 $R$（最大阻值 $999.9\UO$），电阻 $R_{0}$（阻值 $3.0\UO$），电阻 $R_{1}$（阻值 $3.0\UO$），电流表 A（量程为 $200\UmA$，内阻为 $R_{A}=6.0\UO$），开关 S。

实验步骤如下：
\begin{steps}
	\item
	将电阻箱阻值调到最大，闭合开关 S；
	\item
	多次调节电阻箱，记下电流表的示数 $I$ 和电阻箱相应的阻值 $R$；
	\item
	以 $\frac{1}{I}$ 为纵坐标，$R$ 为横坐标，作 $\frac{1}{I}-R$ 图线（用直线拟合）；
	\item
	求出直线的斜率 $k$ 和在纵轴上的截距 $b$。
\end{steps}

回答下列问题：
\begin{enumerate}
	\item
	分别用 $E$ 和 $r$ 表示电源的电动势和内阻，则 $1/I$ 与 $R$ 的关系式为\_\_\_\_\_\_\_\_\_\_。
	\item
	实验得到的部分数据如下表所示，其中电阻 $R=3.0\UO$ 时电流表的示数如图（b）所示，读出数据，完成下表。答：①\_\_\_\_\_\_\_\_\_，②\_\_\_\_\_\_\_\_\_\_\_。

\begin{center}
\begin{tabular}{|c|c|c|c|c|c|c|c|}
\hline
$R/\UO$ & 1.0 & 2.0 & 3.0 & 4.0 & 5.0 & 6.0 & 7.0 \\ \hline
$I/\mathrm{A}$ & 0.143 & 0.125 & ① & 0.100 & 0.091 & 0.084 & 0.077 \\ \hline
$I^{-1}/\mathrm{A}^{-1}$ & 6.99 & 8.00 & ② & 10.0 & 11.0 & 11.9 & 13.0 \\ \hline
\end{tabular}
\end{center}

	\item
	在图（c）的坐标纸上将所缺数据点补充完整并作图，根据图线求得斜率 $k=$\_\_\_\_\_\_$\mathrm{A}^{-1}\Omega^{-1}$，截距 $b=$\_\_\_\_\_\_$\mathrm{A}^{-1}$。
	\item
	根据图线求得电源电动势 $E=$\_\_\_\_\_\_$\mathrm{V}$，内阻 $r=$\_\_\_\_\_\_\_$\UO$。
\end{enumerate}
\threepicture[labelA=2014全国课标Ⅰ23a, labelB=2014全国课标Ⅰ23b, labelC=2014全国课标Ⅰ23c, hA=3.2cm, hB=3.2cm, hC=5.5cm, align=right]
{figs/23a.png}
{figs/23b.png}
{figs/23c.png}
%% answer:
\jdanswer{
\begin{enumerate}
	\item
	$\frac{1}{I}=\frac{3}{E}R+\frac{15+3r}{E}$

	\item
	①$0.110\UA$，②$9.09\mathrm{A}^{-1}$

	\item
	见图（c）答

	\item
	$E=3\UV$，$r=1.0\UO$

\end{enumerate}
}
%% memo:
\memoanswer{
（1）由于电流表 A 与电阻 $R_{1}$ 并联，且 $R_{A}=6.0\UO$，$R_{1}=3.0\UO$，当电流表示数为 $I$ 时，电路总电流应为 $3I$。根据闭合电路欧姆定律有：$E=3I\left(\frac{R_{A}R_{1}}{R_{A}+R_{1}}+R_{0}+R+r\right)$，代入数据整理得：$\frac{1}{I}=\frac{3}{E}R+\frac{15+3r}{E}$。

（2）根据图（b）可得电流表示数为 $110\UmA=0.110\UA$；所以 $1/I=9.09$。

（3）见图（c）答。

（4）根据和图象可得：$\frac{3}{E}=k=1$，$\frac{15+3r}{E}=b=6$，解得：$E=3\UV$，$r=1.0\UO$。
}

\item
%% number: 24
%% paperName: 2014年全国课标Ⅰ第 24 题 12 分
%% typeId: 6
%% type: 计算题
%% chapter: 运动和力
%% point: 牛顿定律的应用
%% method: 无
%% score: 12
%% degree: 650
%% duplicateId: 0
%% body:
公路上行驶的两汽车之间应保持一定的安全距离，当前车突然停止时，后车司机可以采取刹车措施，使汽车在安全距离内停下而不会与前车相碰。通常情况下，人的反应时间和汽车系统的反应时间之和为 $1\Us$，当汽车在晴天干燥沥青路面上以 $108\Ukmh$ 的速度匀速行驶时，安全距离为 $120\Um$。设雨天时汽车轮胎与沥青路面间的动摩擦因数为晴天时的 $\frac{2}{5}$，若要求安全距离仍为 $120\Um$，求汽车在雨天安全行驶的最大速度。
%% answer:
\jdanswer{
$20\Ums$（$72\Ukmh$）
}
%% memo:
\memoanswer{
设路面干燥时，汽车与地面的动摩擦因数为 $\mu_{0}$，刹车时汽车的加速度大小为 $a_{0}$，安全距离为 $s$，反应时间为 $t_{0}$，由牛顿第二定律和运动学公式得

$\mu_{0}mg=ma_{0}$ ①

$s=v_{0}t_{0}+\frac{v_{0}^{2}}{2a_{0}}$ ②

式中，$m$ 和 $v_{0}$ 分别为汽车的质量和刹车前的速度。

设在雨天行驶时，汽车与地面的动摩擦因数为 $\mu$，依题意有

$\mu=\frac{2}{5}\mu_{0}$ ③

设在雨天行驶时汽车刹车的加速度大小为 $a$，安全行驶的最大速度为 $v$，由牛顿第二定律和运动学公式得

$\mu mg=ma$ ④

$s=v t_{0}+\frac{v^{2}}{2a}$ ⑤

联立①②③④⑤式并代入题给数据得

$v=20\Ums$（$72\Ukmh$） ⑥
}

\item
%% number: 25
%% paperName: 2014年全国课标Ⅰ第 25 题 20 分
%% typeId: 6
%% type: 计算题
%% chapter: 电场
%% point: 带电粒子的偏转
%% method: 无
%% score: 20
%% degree: 450
%% duplicateId: 0
%% body:
如图，O、A、B 为同一竖直平面内的三个点，OB 沿竖直方向，$\angle BOA=60^{\circ}$，$OB=\frac{3}{2}OA$。将一质量为 $m$ 的小球以一定的初动能自 O 点水平向右抛出，小球在运动过程中恰好通过 A 点。使此小球带电，电荷量为 $q$（$q>0$），同时加一匀强电场，场强方向与 $\triangle OAB$ 所在平面平行。现从 O 点以同样的初速度沿某一方向抛出此带电小球，该小球通过了 A 点，到达 A 点时的动能是初动能的 3 倍；若该小球从 O 点以同样的初动能沿另一方向抛出，恰好通过 B 点，且到达 B 点时的动能为初动能的 6 倍。重力加速度大小为 $g$。求
\begin{enumerate}
	\item
	无电场时，小球到达 A 点时的动能与初动能的比值；
	\item
	电场强度的大小和方向。
\end{enumerate}
\onepicture[align=right, w=2.5cm]{figs/25.png}
%% answer:
\jdanswer{
\begin{enumerate}
	\item
	$\frac{7}{3}$

	\item
	$\frac{{\sqrt 3 mg}}{{6q}}$，与竖直向下的方向的夹角为 $30^{\circ}$

\end{enumerate}
}
%% memo:
\memoanswer{
（1）设小球的初速度为 $v_{0}$，初动能为 $E_{k0}$，从 O 点运动到 A 点的时间为 $t$，令 $OA=d$，则 $OB=\frac{3}{2}d$，根据平抛运动的规律有

$d\sin60^{\circ}=v_{0}t$ ①

$d\cos60^{\circ}=\frac{1}{2}gt^{2}$ ②

又有 $E_{k0}=\frac{1}{2}mv_{0}^{2}$ ③

由①②③式得

$E_{k0}=\frac{3}{8}mgd$ ④

设小球到达 A 点时的动能为 $E_{kA}$，则

$E_{kA}=E_{k0}+\frac{1}{2}mgd$ ⑤

由④⑤式得

$\frac{E_{kA}}{E_{k0}}=\frac{7}{3}$ ⑥

（2）加电场后，小球从 O 点到 A 点和 B 点，高度分别降低了 $\frac{d}{2}$ 和 $\frac{3d}{2}$，设电势能分别减小 $\Delta E_{pA}$ 和 $\Delta E_{pB}$，由能量守恒及④式得

$\Delta E_{pA}=3E_{k0}-E_{k0}-\frac{1}{2}mgd=\frac{2}{3}E_{k0}$ ⑦

$\Delta E_{pB}=6E_{k0}-E_{k0}-\frac{3}{2}mgd=E_{k0}$ ⑧

在匀强电场中，沿任一直线，电势的降落是均匀的。设直线 OB 上的 M 点与 A 点等电势，M 与 O 点的距离为 $x$，如图，则有

\onepicture[align=right, w=4.0cm]{figs/25_an.png}

$\frac{x}{\frac{3}{2}d}=\frac{\Delta E_{pA}}{\Delta E_{pB}}$ ⑨

解得 $x=d$。MA 为等势线，电场必与其垂线 OC 方向平行。设电场方向与竖直向下的方向夹角为 $\alpha$，由几何关系可得

$\alpha=30^{\circ}$ ⑩

即电场方向与竖直向下的方向的夹角为 $30^{\circ}$。设场强的大小为 $E$，有

$qEd\cos30^{\circ}=\Delta E_{pA}$ \ccnum{⑪}

由④⑦\ccnum{⑪}式得

$E=\frac{\sqrt{3}mg}{6q}$ \ccnum{⑫}
}

\item
%% number: 33
%% paperName: 2014年全国课标Ⅰ第 33 题 9 分
%% typeId: 6
%% type: 计算题
%% chapter: 气体性质
%% point: 气体多过程
%% method: 无
%% score: 9
%% degree: 650
%% duplicateId: 0
%% body:
\textbf{（1）}（6 分）一定量的理想气体从状态 a 开始，经历三个过程 ab、bc、ca 回到原状态，其 $p-T$ 图像如图所示。下列判断正确的是\_\_\_\_\_\_\_\_\_\_\_。
\onepicture[align=right, w=3.5cm]{figs/33.png}
\fivechoices[answer=ADE]
{过程 ab 中气体一定吸热}
{过程 bc 中气体既不吸热也不放热}
{过程 ca 中外界对气体所做的功等于气体所放的热}
{a、b 和 c 三个状态中，状态 a 分子的平均动能最小}
{b 和 c 两个状态中，容器壁单位面积单位时间内受到气体分子撞击的次数不同}

\textbf{（2）}（9 分）一定质量的理想气体被活塞封闭在竖直放置的圆柱形气缸内，气缸壁导热良好，活塞可沿气缸壁无摩擦地滑动。开始时气体压强为 $p$，活塞下表面相对于气缸底部的高度为 $h$，外界的温度为 $T_{0}$。现取质量为 $m$ 的沙子缓慢地倒在活塞的上表面，沙子倒完时，活塞下降了 $h/4$。若此后外界的温度变为 $T$，求重新达到平衡后气体的体积。已知外界大气的压强始终保持不变，重力加速度大小为 $g$。
%% answer:
\jdanswer{
$\frac{9mghT}{4pT_{0}}$
}
%% memo:
\memoanswer{
（1）从 a 到 b 的过程，根据图线过原点可得 $\frac{P_{a}}{T_{a}}=\frac{P_{b}}{T_{b}}$，所以为等容变化过程，气体没有对外做功，外界也没有对气体做功，所以温度升高，只能是吸热的结果，A 正确；从 b 到 c 的过程温度不变，可是压强变小，说明体积膨胀，对外做功，理应内能减少温度降低，而温度不变说明从外界吸热，B 错误；从 c 到 a 的过程，压强不变，根据温度降低说明内能减少，根据改变内能的两种方式及做功和热传递的结果是内能减少，所以外界对气体做的功小于气体放出的热量，C 错误。分子的平均动能与温度有关，状态 a 的温度最低，所以分子平均动能最小，D 正确；b 和 c 两个状态，温度相同，即分子运动的平均速率相等，单个分子对容器壁的平均撞击力相等，根据 b 压强大，可判断状态 b 单位时间内容器壁受到分子撞击的次数多，E 正确。

（2）设气缸的横截面积为 $S$，沙子倒在活塞上后，对气体产生的压强为 $\Delta p$，由玻意耳定律得

$phS=(p+\Delta p)(h-\frac{1}{4}h)S$ ①

解得 $\Delta p=\frac{1}{3}p$ ②

外界的温度变为 $T$ 后，设活塞距底面的高度为 $h'$。根据盖—吕萨克定律，得

$\frac{(h-\frac{1}{4}h)S}{T_{0}}=\frac{h'S}{T}$ ③

解得 $h'=\frac{3T}{4T_{0}}h$ ④

据题意可得 $\Delta p=\frac{mg}{S}$ ⑤

气体最后的体积为 $V=Sh'$ ⑥

联立②④⑤⑥式得

$V=\frac{9mghT}{4pT_{0}}$ ⑦
}

\item
%% number: 34
%% paperName: 2014年全国课标Ⅰ第 34 题 9 分
%% typeId: 6
%% type: 计算题
%% chapter: 光学
%% point: 光的折射
%% method: 无
%% score: 9
%% degree: 450
%% duplicateId: 0
%% body:
\textbf{（1）}（6 分）图（a）为一列简谐横波在 $t=2\Us$ 时波形图，图（b）为媒质中平衡位置在 $x=1.5\Um$ 处的质点的振动图像。P 是平衡位置为 $x=2\Um$ 的质点。下列说法正确的是\_\_\_\_\_\_\_\_。
\twopicture[labelA=2014全国课标Ⅰ34a, labelB=2014全国课标Ⅰ34b, w=5.2cm]
{figs/34a.png}
{figs/34b.png}
\fivechoices[answer=ACE]
{波速为 $0.5\Ums$}
{波的传播方向向右}
{$0\sim2\Us$ 时间内，P 运动的路程为 $8\Ucm$}
{$0\sim2\Us$ 时间内，P 向 $y$ 轴正方向运动}
{当 $t=7\Us$ 时，P 恰好回到平衡位置}

\textbf{（2）}（9 分）一个半圆柱形玻璃砖，其横截面是半径为 $R$ 的半圆，AB 为半圆的直径，O 为圆心，如图所示。玻璃的折射率为 $n=\sqrt 2 $。
\begin{enumerate}
	\item
	一束平行光垂直射向玻璃砖的下表面，若光线到达上表面后，都能从该表面射出，则入射光束在 AB 上的最大宽度为多少？
	\item
	一细束光线在 O 点左侧与 O 点相距 $\frac{\sqrt 3}{2}R$ 处垂直于 AB 从下方入射，求此光线从玻璃砖射出点的位置。
\end{enumerate}
\onepicture[align=right, w=4.0cm]{figs/34c.jpg}
%% answer:
\jdanswer{
\begin{enumerate}
	\item
	$\sqrt 2 R$

	\item
	$\frac{{\sqrt 3 }}{2}R$

\end{enumerate}
}
%% memo:
\memoanswer{
（1）根据图（a）的波形图判断机械波的波长 $\lambda=2\Um$，根据图（b）可得振动周期 $T=4\Us$，所以波速 $v=\lambda/T=0.5\Ums$，A 正确；根据图（b）可判断 $x=1.5\Um$ 的质点在 $t=2\Us$ 振动方向为 $y$ 轴负方向，在图（a）中，根据质点振动方向和传播方向在图像同一侧可判断波的传播方向向左，B 错误；$t=2\Us$ 时质点 P 在最低点，根据周期 $T=4\Us$，可知 $t=0$ 时质点 P 在最高点，所以 $0\sim2\Us$ 时间内质点 P 通过的路程为 2 倍的振幅即 $8\Ucm$，C 正确；$0\sim2\Us$ 质点 P 向 $y$ 轴负方向运动，D 错误；$t=2\Us$ 到 $t=7\Us$ 共经 $5\Us$ 为 $5T/4$，所以质点 P 刚好回到平衡位置，E 正确。

（2）（i）在 O 点左侧，设从 E 点射入的光线进入玻璃砖后在上表面的入射角恰好等于全反射的临界角 $\theta$，则 OE 区域的入射光线经上表面折射后都能从玻璃砖射出，如图。由全反射条件有

$\sin\theta=\frac{1}{n}$ ①

由几何关系有

$OE=R\sin\theta$ ②

由对称性可知，若光线都能从上表面射出，光束的宽度最大为

$l=2OE$ ③

\onepicture[align=right, w=4.5cm]{figs/34_an1.jpg}

联立①②③式，代入已知数据得

$l=\sqrt{2}R$ ④

（ii）设光线在距 O 点 $\frac{\sqrt{3}}{2}R$ 的 C 点射入后，在上表面的入射角为 $\alpha$，由几何关系及①式和已知条件得

$\alpha=60^{\circ}>\theta$ ⑤

光线在玻璃砖内会发生三次全反射，最后由 G 点射出，如图。由反射定律和几何关系得

\onepicture[align=right, w=4.5cm]{figs/34_an2.jpg}

$OG=OC=\frac{\sqrt{3}}{2}R$ ⑥

射到 G 点的光有一部分被反射，沿原路返回到达 C 点射出。
}

\item
%% number: 35
%% paperName: 2014年全国课标Ⅰ第 35 题 9 分
%% typeId: 6
%% type: 计算题
%% chapter: 运动和力
%% point: 动量守恒定律
%% method: 无
%% score: 9
%% degree: 450
%% duplicateId: 0
%% body:
\textbf{（1）}（6 分）关于天然放射性，下列说法正确的是\_\_\_\_\_\_\_\_\_\_\_\_。
\fivechoices[answer=BCD]
{所有元素都可能发生衰变}
{放射性元素的半衰期与外界的温度无关}
{放射性元素与别的元素形成化合物时仍具有放射性}
{$\alpha$、$\beta$ 和 $\gamma$ 三种射线中，$\gamma$ 射线的穿透能力最强}
{一个原子核在一次衰变中可同时放出 $\alpha$、$\beta$ 和 $\gamma$ 三种射线}

\textbf{（2）}（9 分）如图，质量分别为 $m_{A}$、$m_{B}$ 的两个弹性小球 A、B 静止在地面上方，B 球距地面的高度 $h=0.8\Um$，A 球在 B 球的正上方。先将 B 球释放，经过一段时间后再将 A 球释放。当 A 球下落 $t=0.3\Us$ 时，刚好与 B 球在地面上方的 P 点处相碰，碰撞时间极短，碰后瞬间 A 球的速度恰为零。已知 $m_{B}=3m_{A}$，重力加速度大小 $g=10\Umsq$。忽略空气阻力及碰撞中的动能损失。求：
\begin{enumerate}
	\item
	B 球第一次到达地面时的速度；
	\item
	P 点距离地面的高度。
\end{enumerate}
\onepicture[align=right, w=3.5cm]{figs/35.png}
%% answer:
\jdanswer{
\begin{enumerate}
	\item
	$4\Ums$

	\item
	$0.75\Um$

\end{enumerate}
}
%% memo:
\memoanswer{
（1）只有原子序数超过 83 的元素才都能发生衰变，A 错误；放射性元素的半衰期决定于由原子核内部的结构，与外界温度及化学作用等无关，B 正确；放射性元素其放射性来自于原子核内部的，与其他元素形成化合物并没有改变其内部原子核结构所以仍具有放射性，C 正确；$\alpha$、$\beta$ 和 $\gamma$ 三种射线中，$\gamma$ 射线能量最高，穿透能力最强，D 正确；一个原子核在一次衰变中，要是 $\alpha$ 衰变、要么是 $\beta$ 衰变，同时伴随着能量的释放，即 $\gamma$ 射线，E 错误。

（2）（i）设 B 球第一次到达地面时的速度大小为 $v_{B}$，由运动学公式有

$v_{B}=\sqrt{2gh}$ ①

将 $h=0.8\Um$ 代入上式，得

$v_{B}=4\Ums$ ②

（ii）设两球相碰前后，A 球的速度大小分别为 $v_{1}$ 和 $v_{1}'$（$v_{1}'=0$），B 球的速度分别为 $v_{2}$ 和 $v_{2}'$，由运动学规律可得

$v_{1}=gt$ ③

由于碰撞时间极短，重力的作用可以忽略，两球相碰前后的动量守恒，总动能保持不变。规定向下的方向为正，有

$m_{A}v_{1}+m_{B}v_{2}=m_{B}v_{2}'$ ④

$\frac{1}{2}m_{A}v_{1}^{2}+\frac{1}{2}m_{B}v_{2}^{2}=\frac{1}{2}m_{B}v_{2}'^{2}$ ⑤

设 B 球与地面相碰后的速度大小为 $v_{B}'$，由运动学及碰撞的规律可得

$v_{B}'=v_{B}$ ⑥

设 P 点距地面的高度为 $h'$，由运动学规律可得

$h'=\frac{v_{B}'^{2}-v_{2}^{2}}{2g}$ ⑦

联立②③④⑤⑥⑦式，并代入已知条件可得

$h'=0.75\Um$ ⑧
}

\end{enumerate}

\end{document}
